% Generated draft: requires mathematical and source review.
\documentclass[11pt]{article}
\usepackage[margin=25mm]{geometry}
\usepackage{fontspec}
\setmainfont{Latin Modern Roman}
\usepackage{amsmath,amssymb,mathtools,unicode-math}
\setmathfont{Latin Modern Math}
\usepackage{enumitem,needspace,xurl,hyperref}
\setlength{\parindent}{0pt}
\setlength{\parskip}{5pt}
\setlength{\emergencystretch}{3em}
\newcommand{\N}{\mathbb{N}}
\newcommand{\Z}{\mathbb{Z}}
\newcommand{\Q}{\mathbb{Q}}
\newcommand{\R}{\mathbb{R}}
\newcommand{\C}{\mathbb{C}}
\newcommand{\F}{\mathbb{F}}
\newcommand{\A}{\mathbb{A}}
\newcommand{\PP}{\mathbb{P}}
\newcommand{\E}{\mathbb{E}}
\newcommand{\eps}{\varepsilon}
\newcommand{\dd}{\,\mathrm{d}}
\begin{document}
\section*{UnsolvedMath ID: 30004370}
\textbf{Source problem code:} catalogue entry
\subsection{Definitions and mathematical statement}
The supplied record attributes the following realization question to Brown--Wemyss: does every finite-dimensional Jacobi algebra arise as the contraction algebra of a flopping curve? The record does not specify the ground field, the convention for Jacobi algebras, or the permitted geometric contractions. These qualifications remain unresolved.

Here are explicit algebraic conventions that distinguish some of the missing choices. Let \(k\) be a field and \(Q\) a finite quiver. Its path algebra \(kQ\) has the directed paths as a basis, with multiplication by concatenation when possible and zero otherwise. For a potential \(W\), meaning a linear combination of cyclic paths considered modulo cyclic rotation, the cyclic derivative with respect to an arrow \(a\) is determined by
\[
\partial_a(a_1\cdots a_m)
 =\sum_{\substack{1\leq i\leq m\\a_i=a}}
 a_{i+1}\cdots a_m a_1\cdots a_{i-1}.
\]
The ordinary Jacobi algebra is
\[
J_k(Q,W)=kQ/\langle\partial_aW:a\in Q_1\rangle,
\]
where \(Q_1\) denotes the arrow set. Alternatively, writing \(\mathfrak m\) for the ideal generated by arrows, one can use the completed path algebra
\[
\widehat{kQ}=\varprojlim_{n\geq1}kQ/\mathfrak m^n
\]
and a formal potential to obtain
\[
\widehat J_k(Q,W)=\widehat{kQ}/\overline{\langle\partial_aW:a\in Q_1\rangle}.
\]
The bar denotes closure in the arrow-adic topology. Finite dimensionality means finite dimension as a \(k\)-vector space. The supplied record does not select between these conventions or specify restrictions on \(Q\) or \(W\).

In a local module presentation of a contraction algebra, the algebraic construction takes the form
\[
A_{\mathrm{con}}=\operatorname{End}_R(M)/[R],
\]
where \(R\) is the local base ring, \(M\) is a module supplied by the geometric construction, and \([R]\) is the ideal of endomorphisms that factor through a finite direct sum of copies of \(R\). This description does not specify which geometric contractions or modules are admissible; those hypotheses must be recovered from the cited source. In particular, the record does not define the ambient dimension, singularity conditions, or whether the exceptional locus must be one irreducible curve.

With these unresolved choices displayed, the supplied conjecture has the template
\[
\forall A\in\mathcal J\quad
\exists f\in\mathcal F\quad
A\cong A_{\mathrm{con}}(f),
\]
where \(\mathcal J\) is the intended class of finite-dimensional Jacobi algebras and \(\mathcal F\) is the intended class of flopping contractions. Neither class is fully specified by the record. The intended meaning of realization, including algebra isomorphism versus Morita equivalence, also requires confirmation. This is consequently a statement-only draft of the supplied question, with its missing hypotheses exposed, rather than a fully determined universal assertion.
\subsection{Short English statement}
The Brown--Wemyss question asks whether every finite-dimensional Jacobi algebra in the intended class can be realized geometrically as the contraction algebra of a flopping curve. The supplied record labels the problem open, but leaves the algebraic and geometric scope insufficiently specified. Its dated assessment describes nonnegative refined invariants as evidence, not a proof.
\subsection{Sources}
\begin{itemize}
\item Local registry: \texttt{data/unsolvedmath/problems.json}, permanent UnsolvedMath ID \(30004370\), source problem code \texttt{catalogue entry}, title \emph{Jacobi Algebras as Contraction Algebras}. The original statement attributes the conjecture to Brown--Wemyss; the separate proposer field is empty, and the supplied proposed year is 2020.
\item Supplied report citation: \emph{Representation Theory of Quivers and Finite Dimensional Algebras} (2021), Oberwolfach Reports, pp.~143--230, DOI \texttt{10.4171/owr/2020/3}. \url{https://doi.org/10.4171/owr/2020/3}. The background describes Oberwolfach Reports as documenting mathematical workshops and their open-problem discussions. No exact page or passage for this conjecture is supplied.
\item Davison, \emph{Refined invariants of finite-dimensional Jacobi algebras}. \url{https://arxiv.org/abs/1903.00659}.
\item Brown--Wemyss, \emph{Noncommutative singularity theory}. \url{https://arxiv.org/abs/2410.21500}.
\item Supplied literature assessment dated exactly \texttt{2026-08-22}: the Brown--Wemyss realization conjecture for all relevant finite-dimensional Jacobi algebras remains open; nonnegative refined invariants provide evidence but not a proof. This is supplied dated curation derived from the OWR record, not an independently verified literature assessment. The background explicitly requests independent verification of time-sensitive status and cited claims.
\item The registry supplies status \texttt{open} and statement status \texttt{corrected\_verified}, together with a claim that the clean statement was checked against the official report and preserves its hypotheses, notation, and target. These are supplied provenance claims, not verification performed for this draft.
\end{itemize}
\end{document}
